\chapter{Aspiration Pneumonia}
\index{Aspiration Pneumonia}
\label{sec:Aspiration Pneumonia}

\section{Overview}
Aspiration pneumonia is a pulmonary infection caused by the inhalation of oropharyngeal or gastric contents into the lower respiratory tract, leading to chemical pneumonitis and subsequent bacterial infection. It most commonly affects dependent lung segments, particularly the right lower lobe due to the more vertical orientation of the right main bronchus. Aspiration pneumonia is distinct from aspiration pneumonitis (Mendelson syndrome), which is a chemical injury without infection.
\section{Classification}

\begin{description}[font=\normalfont\bfseries,leftmargin=3em,labelwidth=2.5em]
  \item[Cause] Infectious (bacterial) following aspiration of colonized oropharyngeal material
  \item[Organ System] Lungs/Respiratory
  \item[Pathophysiology] Aspiration of oropharyngeal or gastric contents into lower airways causing chemical injury followed by bacterial infection with mixed aerobic and anaerobic flora
  \item[Duration] Acute
  \item[Onset] Sudden (chemical) to Gradual (infectious over days)
  \item[Mode of Transmission] Non-communicable (aspiration of endogenous flora)
  \item[Clinical Course] Variable; can progress rapidly to severe pneumonia or lung abscess
  \item[Severity] Mild to Severe
  \item[Extent of Disease] Localized (dependent lung segments, often right lower lobe)
  \item[Age Group] Older adults, young children, and neurologically impaired patients
  \item[Epidemiology] Accounts for 5--15\% of community-acquired pneumonia cases; high incidence in nursing homes and hospitalized patients with dysphagia
  \item[ICD-11 Code] CA22.0
\end{description}

\section{Causes}
Aspiration pneumonia results from the entry of bacteria-laden oropharyngeal secretions or gastric contents into the lower airways, typically in patients with impaired swallowing, reduced cough reflex, or decreased level of consciousness. The primary pathogens include anaerobic oral flora (Peptostreptococcus, Bacteroides, Prevotella, Fusobacterium) often mixed with aerobic bacteria (Streptococcus pneumoniae, Haemophilus influenzae, Staphylococcus aureus). Hospital-acquired aspiration may involve more resistant organisms such as Pseudomonas aeruginosa and Enterobacteriaceae.

\section{Risk Factors}

\begin{itemize}[noitemsep]
  \item Dysphagia from neurologic disorders (stroke, Parkinson disease, dementia, amyotrophic lateral sclerosis)
  \item Decreased level of consciousness (alcohol or drug intoxication, general anesthesia, sedation, seizure)
  \item Impaired cough reflex due to neurologic injury or medications
  \item Nasogastric or endotracheal tubes compromising laryngeal function
  \item Gastroesophageal reflux disease with large-volume regurgitation
  \item Poor oral hygiene and periodontal disease leading to high bacterial load
  \item Advanced age with frailty and institutionalization
\end{itemize}

\section{Signs and Symptoms}

\subsection{Early symptoms}
\begin{itemize}[noitemsep]
  \item Early manifestations of aspiration pneumonia may be subtle and nonspecific
\end{itemize}

\subsection{Common symptoms}
\begin{itemize}[noitemsep]
  \item \subsection{Common symptoms}
  \item \textbf{Common symptoms:}
  \item Cough productive of foul-smelling or purulent sputum
  \item Fever and rigors
  \item Dyspnea and tachypnea
  \item Pleuritic chest pain
  \item Crackles or bronchial breath sounds in affected lung zones
  \item Hypoxia with decreased oxygen saturation
  \item Confusion or altered mental status (especially in elderly)
  \item Tachycardia and signs of systemic inflammatory response
  \item Pain or discomfort in the affected area
  \item Difficulty performing daily activities
\end{itemize}

\subsection{Less common symptoms}
\begin{itemize}[noitemsep]
  \item Atypical or less common presentations of aspiration pneumonia may occur
\end{itemize}

\subsection{Emergency warning signs}
\begin{itemize}[noitemsep]
  \item Severe or worsening symptoms of aspiration pneumonia require immediate medical evaluation
\end{itemize}
\section{Progression and Stages}
Chemical pneumonitis from acidic gastric contents develops within hours of aspiration, presenting with acute dyspnea and hypoxemia even without infection. Bacterial aspiration pneumonia typically evolves over 1--3 days as aspirated organisms multiply in the lung parenchyma, producing a consolidated infection that appears as dependent-segment opacities on chest imaging. Untreated or inadequately treated infection can progress to necrotizing pneumonia, lung abscess formation, or empyema over the course of 7--14 days.

\section{Complications}

\begin{itemize}[noitemsep]
  \item Lung abscess (cavitary lesion from necrotizing infection)
  \item Empyema (purulent pleural effusion requiring drainage)
  \item Bronchopleural fistula
  \item Bacteremia with sepsis and septic shock
  \item Respiratory failure requiring mechanical ventilation
  \item Reduced quality of life
  \item Emotional and psychological effects
\end{itemize}

\section{Diagnosis}

\begin{itemize}[noitemsep]
  \item Clinical history of aspiration event or presence of risk factors (dysphagia, altered consciousness)
  \item Chest radiograph or CT showing opacities in dependent segments (right lower lobe, superior segments of lower lobes)
  \item Sputum gram stain and culture (aerobic and anaerobic)
  \item Blood cultures to identify bacteremia
  \item Complete blood count with leukocytosis and left shift
  \item Procalcitonin and C-reactive protein as inflammatory markers
  \item Imaging tests to evaluate for complications (cavitation, empyema)
\end{itemize}

\section{Prevention}

\begin{itemize}[noitemsep]
  \item Dysphagia screening and swallowing assessment in at-risk patients (post-stroke, elderly)
  \item Modified food textures and thickened liquids for patients with impaired swallowing
  \item Head-of-bed elevation to 30--45 degrees during and after feeding
  \item Good oral hygiene to reduce bacterial colonization of the oropharynx
  \item Proper tube feeding management with gastric residual volume monitoring
  \item Minimize sedative and antipsychotic medications in susceptible patients
  \item Go for regular check-ups
  \item Follow your doctor's screening recommendations
\end{itemize}

\section{Treatment Options}

\begin{itemize}[noitemsep]
  \item Empiric broad-spectrum antibiotics covering anaerobes and aerobes (beta-lactam/beta-lactamase inhibitor, clindamycin, or carbapenem for severe cases)
  \item Targeted antibiotic therapy based on culture sensitivity results
  \item Supplemental oxygen and respiratory support as needed
  \item Chest physiotherapy and postural drainage to mobilize secretions
  \item Source control: dental treatment for periodontal disease if present
  \item Tube feeding modifications or percutaneous gastrostomy for severe dysphagia
  \item Lifestyle changes and self-care
  \item Surgery when needed (drainage of empyema, lung resection for abscess)
\end{itemize}

\section{Living with It}

\begin{itemize}[noitemsep]
  \item Follow your treatment plan as prescribed
  \item Keep up with regular appointments
  \item Adjust your daily habits as needed
  \item Reach out to family, friends, or support groups
  \item Keep learning about your condition
\end{itemize}

\section{Outlook}
With prompt diagnosis and appropriate antibiotic therapy, most patients with aspiration pneumonia recover fully, though recovery may be prolonged in elderly or debilitated patients. Mortality ranges from 5--20\% depending on the severity of infection, comorbidities, and the patient's functional status.
